\subsubsection{Alimentação}

% Justificar a estratégia de alimentação a partir do cenário de uso (dispositivo de cabeceira fixo vs. portátil), discutindo implicações de autonomia, segurança elétrica e norma aplicável a equipamentos médicos.
% Decisão pendente: dispositivo fixo de cabeceira ou portátil/wearable?

% Se fixo: fonte chaveada 5V/3A com proteção contra sobretensão; mais alinhado ao escopo do TCC.
% Se portátil: bateria Li-Po + TP4056 (carga) + step-up 5V + circuito de proteção; adiciona complexidade de gerenciamento de energia que pode comprometer o cronograma.

% Definir tempo mínimo de autonomia caso a opção portátil seja mantida e qual a estratégia de fallback em caso de queda de energia.

A estratégia de alimentação do sistema foi definida considerando seu uso como um dispositivo fixo de monitoramento à beira do leito. Embora dispositivos alimentados por bateria sejam amplamente utilizados em aplicações portáteis, sistemas destinados ao monitoramento clínico contínuo demandam elevada disponibilidade e operação ininterrupta. Nesses cenários, a alimentação exclusivamente por baterias pode se tornar impraticável devido à necessidade de recargas periódicas e ao risco de interrupções durante o monitoramento \cite{schweber2022}.

Diante desses requisitos, optou-se por uma arquitetura alimentada diretamente pela rede elétrica. Essa abordagem elimina limitações relacionadas à autonomia energética e reduz a complexidade do projeto ao dispensar circuitos dedicados ao carregamento e gerenciamento de baterias. Além disso, permite maior previsibilidade operacional e simplifica a manutenção do dispositivo.

Para o protótipo, prevê-se a utilização de uma fonte chaveada AC-DC de 5 V e 3 A, responsável por alimentar a plataforma de processamento, os sensores, o display e os demais periféricos do sistema. A adoção de uma fonte comercial simplifica a implementação do hardware e contribui para a segurança operacional do equipamento, estando alinhada às boas práticas de segurança elétrica aplicadas a equipamentos eletromédicos, conforme estabelecido pela norma IEC 60601-1 da \textit{International Electrotechnical Commission} (IEC) \citeonline{iec60601}.

Por fim, a arquitetura de alimentação proposta fornece energia suficiente para a execução simultânea das tarefas de aquisição de sinais fisiológicos, processamento embarcado e inferência dos modelos de aprendizado de máquina. Dessa forma, contribui para a confiabilidade e estabilidade operacional do sistema no cenário de monitoramento contínuo para o qual foi concebido.
